\documentclass[12pt]{article}
\usepackage[utf8]{inputenc}
\usepackage[T1]{fontenc}
\usepackage{hyperref}
\title{Tofin Gbe Language Resources for K-12 Curriculum Development}
\author{Compiled by Claude Sonnet 5 (claude-sonnet-5) for Andr\'as Kornai}
\date{}
\begin{document}
\maketitle
\section{Introduction}
Tofin (T\textopenbullet fin, ISO 639-3 tfi) is a Gbe language of southern Benin, spoken in and around Porto-Novo. Wikipedia gives about 66,000 speakers (2006); Global Recordings gives 90,000. The search found very little that is educational. Nothing found shows Tofin being taught in schools.
\section{Primary Educational Resources}
No primers, textbooks or school programmes for Tofin turned up in search.
\section{Linguistic and Reference Materials}
\begin{itemize}
\item Overview article on the language: \url{https://en.wikipedia.org/wiki/Tofin_language}
\item Gabi Schoch (2011), SIL sociolinguistic survey of the Tofin area, a volume in a series on Gbe communities in Benin and Togo (reference record only): \url{https://glottolog.org/resource/reference/id/161211}
\item Related variety T\textopenbullet li, which Glottolog treats as a Tofin dialect and Ethnologue treats as a Gun dialect: \url{https://en.wikipedia.org/wiki/T%C9%94li_language}
\end{itemize}
\section{Religious and Community Resources}
\begin{itemize}
\item Joshua Project people-group profile: \url{https://www.joshuaproject.net/people_groups/15528}
\item Scripture Earth index entry for tfi, whose contents were not verified: \url{https://www.scriptureearth.org/00i-Scripture_Index.php?iso=tfi}
\end{itemize}
\section{Digital and Audio Resources}
Global Recordings Network lists a page for tfi, whose contents were not verified: \url{https://globalrecordings.net/en/language/tfi}
\section{Curriculum Development Context}
Little material exists. A curriculum would need new orthography work, built on the Gbe literature and on Fon and Gun, which are related and mutually intelligible according to Joshua Project. Dialect boundaries with T\textopenbullet li are disputed.
\end{document}